% main_report78.tex
% SAMBP Technical Report TR-78/2028
% Adaptive Protection Under N-1 Contingency
\documentclass[11pt, a4paper]{article}

\usepackage[T1]{fontenc}
\usepackage[utf8]{inputenc}
\usepackage[margin=25mm]{geometry}
\usepackage{amsmath, amssymb, amsthm}
\usepackage{booktabs, array, multirow, longtable}
\usepackage{graphicx}
\usepackage[table]{xcolor}
\usepackage[hidelinks]{hyperref}
\usepackage{pgfplots}
\pgfplotsset{compat=1.18}
\usepackage{tikz}
\usetikzlibrary{arrows.meta, positioning, calc, fit, backgrounds,
                shapes.geometric, shapes.misc}
\usepackage{caption}
\usepackage{subcaption}
\usepackage{parskip}
\usepackage{siunitx}
\usepackage{pifont}
\usepackage{cite}

\definecolor{sambpblue}{RGB}{0,70,140}
\definecolor{okgreen}{RGB}{0,130,60}
\definecolor{alertred}{RGB}{180,0,0}
\definecolor{warnoran}{RGB}{200,100,0}

\newcommand{\pu}{\,\text{pu}}
\newcommand{\ms}{\,\text{ms}}
\newcommand{\cmark}{\ding{51}}
\newcommand{\xmark}{\ding{55}}

\newtheorem{finding}{Finding}
\newtheorem{remark}{Remark}

\title{\textbf{\textsf{\color{sambpblue}SAMBP Technical Report TR-78/2028}}\\[6pt]
  \Large Adaptive Protection Under N-1 Contingency\\[4pt]
  \normalsize Zone Reach Recalculation, Topology Change via GOOSE,
  Blinder Adjustment: 100-Scenario IEEE~39-Bus Study}
\author{SAMBP Research Group\\[2pt]
  \small Smart Adaptive Multi-Bus Protection Research, IIT Madras}
\date{April 2028\quad\textbar\quad Chapter~5 (Distance Protection)\quad\textbar\quad
  Prerequisite: TR-72 (Wide-Area PMU)}

\begin{document}
\maketitle
\tableofcontents
\vspace{6pt}\hrule\vspace{10pt}

\section{Background and Objectives}
\label{sec:background}

\subsection{Motivation}

Distance relay settings are traditionally calculated for the base-case (N-0)
network topology. Under N-1 contingency (loss of a line), the network impedances
change, and fixed Zone~1/2 settings may over-reach or under-reach, causing
miscoordination or missed trips~\cite{Horowitz2008}. In IBR-rich networks,
N-1 contingencies also change the fault current distribution (infeed changes),
compounding the reach error.

The SAMBP adaptive protection framework addresses this by:
\begin{enumerate}
  \item Detecting topology changes via breaker status GOOSE messages.
  \item Recalculating zone reaches from a pre-computed N-1 look-up table.
  \item Adjusting blinder settings for load encroachment under heavy power transfer.
\end{enumerate}

\subsection{Objectives}

\begin{enumerate}
  \item Implement topology change detection from IEC~61850 GOOSE breaker status.
  \item Build an N-1 zone reach table for IEEE~39-bus (39 line outage contingencies).
  \item Develop blinder adjustment rules for load encroachment.
  \item Validate 100 contingency scenarios: misoperation rate $< 2\%$.
\end{enumerate}

\section{Theory and Methodology}
\label{sec:theory}

\subsection{Topology Change Detection}

Each circuit breaker publishes its status (OPEN/CLOSE) on the GOOSE bus
at $\leq 4\,\ms$ latency. The topology change detection algorithm:
\begin{enumerate}
  \item Monitors the GOOSE dataset for breaker status changes.
  \item On detecting a change, updates the network topology model
    (removes the outaged element from $\mathbf{Y}_{\mathrm{bus}}$).
  \item Triggers a zone reach recalculation from the N-1 look-up table.
  \item Applies new relay settings within $t_{\mathrm{adapt}} \leq 50\,\ms$
    (GOOSE detection: $\leq 4\,\ms$ + table lookup: $< 1\,\ms$ +
    relay setting write: $\leq 45\,\ms$).
\end{enumerate}

\subsection{N-1 Zone Reach Table}

For each N-1 contingency $c$ (line outage), the Zone~1 and Zone~2 reaches
for relay $r$ are pre-computed:

\begin{align}
  Z_1^{(c)} &= 0.80 \, Z_{\mathrm{line}}^{(c)} / (1 + \kibr^{(c)})
  \label{eq:z1}\\
  Z_2^{(c)} &= 1.20 \, Z_{\mathrm{line}}^{(c)} + 0.50 \, Z_{\mathrm{adj}}^{(c)}
  \label{eq:z2}
\end{align}

\noindent where $Z_{\mathrm{line}}^{(c)}$ is the protected line impedance
(unchanged by the outage, but the source impedance at both ends changes),
$\kibr^{(c)}$ is the IBR current ratio under contingency $c$ (from TR-77 SE),
and $Z_{\mathrm{adj}}^{(c)}$ is the adjacent line Zone~1 reach under the contingency.

The table is pre-computed offline for all 39 N-1 contingencies and stored
in the SAMBP relay database (total size: $39 \times N_{\mathrm{relays}} \times 4$
parameters).

\subsection{Blinder Adjustment for Load Encroachment}

Under heavy power transfer (post N-1), the load impedance locus may enter the
Zone~2 characteristic. The blinder adjustment restricts the angular reach:

\begin{equation}
  \phi_{\mathrm{blinder}}^{(c)} = \phi_{\mathrm{base}} -
    \Delta\phi \cdot \frac{P_{\mathrm{transfer}}^{(c)}}{P_{\mathrm{rated}}}
  \label{eq:blinder}
\end{equation}

\noindent where $\phi_{\mathrm{base}} = 30^\circ$, $\Delta\phi = 10^\circ$,
and $P_{\mathrm{transfer}}^{(c)}/P_{\mathrm{rated}}$ is the normalised
post-contingency transfer level from the SE.

\section{Simulation Results}
\label{sec:results}

\subsection{100-Scenario Results}

\begin{table}[ht]
\centering
\caption{TR-78 100-scenario results on IEEE~39-bus.}
\label{tab:results}
\begin{tabular}{lccccc}
\toprule
Contingency type & Scenarios & Correct & Misop.\ & Misop.\ rate & Status \\
\midrule
Single line outage (light load)  & 40 & 40 & 0 & 0\%   & \cmark \\
Single line outage (heavy load)  & 30 & 29 & 1 & 3.3\% & \xmark \\
Double contingency (N-2)         & 20 & 20 & 0 & 0\%   & \cmark \\
IBR source outage (N-1 IBR)      & 10 & 10 & 0 & 0\%   & \cmark \\
\midrule
\textbf{Total} & \textbf{100} & \textbf{99} & \textbf{1}
  & \textbf{1.0\%} & \textbf{\color{okgreen}PASS ($< 2\%$)} \\
\bottomrule
\end{tabular}
\end{table}

\begin{remark}
The one misoperation is a heavy-load scenario where the blinder adjustment
is insufficient: the blinder angle $(20^\circ)$ still admits a load impedance
locus entry for $\cos\phi = 0.85$ lagging at 95\% transfer. Resolution:
reduce $\Delta\phi$ to $15^\circ$ for extreme load angles, applied in v1.1
of the N-1 table.
\end{remark}

\subsection{Adaptation Time Performance}

\begin{table}[ht]
\centering
\caption{Zone reach adaptation time after N-1 event.}
\label{tab:adapt_time}
\begin{tabular}{lc}
\toprule
Component & Time (ms) \\
\midrule
GOOSE breaker status detection & $\leq 4$ \\
N-1 table lookup & $< 1$ \\
New reach parameter write & $\leq 45$ \\
\midrule
Total adaptation time & $\leq 50$ \\
\midrule
Zone~2 time delay (typical) & 400 \\
\bottomrule
\end{tabular}
\end{table}

The 50\,ms adaptation time is $8\times$ faster than the Zone~2 time delay
(400\,ms typical), ensuring that the correct Zone~2 reach is active before
any Zone~2 trip decision under the N-1 topology.

\section{Conclusion}
\label{sec:conclusion}

\begin{finding}[N-1 adaptive protection: 1\% misoperation rate]
The N-1 adaptive zone reach framework achieves 99/100 = 99\% correct
operation (1\% misoperation), well below the $< 2\%$ criterion.
The adaptation time of $\leq 50\,\ms$ ensures correct settings before
any Zone~2 decision under the new topology.
\end{finding}

\begin{finding}[GOOSE topology detection enables adaptive settings]
The IEC~61850 GOOSE breaker status detection provides topology change
notification within 4\,ms, enabling the SAMBP framework to recalculate
zone reaches 50\,ms after the contingency --- a fundamental improvement
over manual periodic setting group changes.
\end{finding}

\begin{thebibliography}{99}
\bibitem{Horowitz2008}
S.~H.~Horowitz and A.~G.~Phadke,
\emph{Power System Relaying}, 3rd ed., Wiley, Chichester, 2008.

\bibitem{Erickson2010}
J.~Erickson \emph{et al.},
``Adaptive relaying concepts to improve transmission systems protection,''
in \emph{Proc.\ IEEE PES T\&D Conference}, New Orleans, LA, 2010.

\bibitem{Yalla1992}
M.~V.~V.~S.~Yalla,
``A digital multifunction protective relay,''
\emph{IEEE Trans.\ Power Deliv.}, vol.~7, no.~1, pp.~193--201, 1992.

\bibitem{IEC61850GOOSE}
IEC~61850-9-2:2011, \emph{Communication Networks and Systems for Power
Utility Automation --- Part~9-2: Sampled Values over ISO/IEC~8802-3},
IEC, Geneva, 2011.

\bibitem{Venkataraman2014}
S.~Venkataraman \emph{et al.},
``Adaptive distance relay settings for transmission systems with high
penetration of DG,''
\emph{Electr.\ Power Syst.\ Res.}, vol.~134, pp.~44--53, 2016.
\end{thebibliography}

\end{document}
